\documentclass{subfiles}
\begin{document}
    Consider the case in which the keys are choosed in such a way that, after the hashing,
        the average retrival time is \(\ofT{n}\), a situation that can happen with a non-zero 
        probability. The use of fixed hash functions may lead to such situations.
        A simple solution, is that of choosing the hash function at random from a family of well-designed 
        hash functions.

    Here, we first show why this kind of solution works, 
        then, we show how to one of such families is defined.

    Let \(\H\) be a finite collection of hash functions. 
        We say that \(\H\) is universal if, for any pair of keys \(k, l \in K\), 
        the number of hash function such that \(h(k) = h(l)\) is no more than \(\sfrac{\abs{\H}}{m}\).
        In other words, \(\H\) is universal if the probability of a collision is less than \(\sfrac{1}{m}\).

    The following theorem shows that universal has functions give a good-average case.
    \begin{theorem}
        Let \(h\) be a function that has been used to hash \(n\) keys into the table \(T\) of size \(m\).
            If key \(k\) is not in the table, then the expected lenght of the list \(k\) hashes to is at most
            \(\sfrac{n}{m} = \alpha\). Viceversa, if key \(k\) is in the table, 
            the expected lenght of the list storing \(k\) is at most \(1 + \alpha\).
    \end{theorem}
    \begin{proof*}
        Let \(X_{kl}\) be the indicator random variable \(I\set{h(k) = h(l)}\).
            By definition of universal universal collection of hash functions,
            \(\Prob{h(k) = h(l)} \le \sfrac{1}{m}\). By a simple result in probability theory,
            we know \(\E[X_{kl}]{} \le \sfrac{1}{m}\).

        Let \(Y_{k}\) denote the number of keys \(l \ne k\) that maps to the same slot as \(k\).
        We have
        \[\begin{aligned}
            \E[Y_{k}]{} & = \E[\sum_{l \in K, l \ne k}]{} \\ 
                     & = \sum_{l \in K, l \ne k} \E[X_{kl}]{} \\ 
                     & \le \sum_{l \in K, l \ne k} \frac{1}{m}                  
        \end{aligned}\]
        
        If \(n_{h(k)}\) is the lenght of the list \(k\) hashes to, then 
        \begin{itemize}
            \item if \(k \notin T\), then \(\E[n_{h(k)}]{} = \E[Y_{k}]{}\) and 
                \(\abs{\set{l}[l \in T \land l \ne k]} = n\).
                Therefore, \(\E[n_{h(k)}]{} \le \sfrac{n}{m} = \alpha\).

            \item if \(k \in T\), then \(\E[n_{h(k)}]{} = \E[Y_{k}]{} + 1\) and
                \(\abs{\set{l}[l \in T \land l \ne k]} = n - 1\).
                Hence, \(\E[n_{h(k)}]{} \le (n - 1) / m  + 1 = 1 + \alpha - \sfrac{1}{m}  \le 1 + \alpha\)
        \end{itemize}\qed
    \end{proof*}

    \subsubsection{Designing a universal class of hash functions}
    \subfile{../subsub/3.2.1 - Designing a universal class of hash functions}
\end{document}
